% !Mode:: "TeX:UTF-8"
% chapters/cha.2_主体内容.tex — 第2章示例：主体内容（含图、表、公式示例）

\chapter{主体内容示例}

\section{公式示例}

行内公式：设损失函数为 $\mathcal{L} = \frac{1}{N}\sum_{i=1}^{N}(y_i - \hat{y}_i)^2$。

行间公式如式~(\ref{eq:loss}) 所示：

\begin{equation}
    \mathcal{L} = \frac{1}{N}\sum_{i=1}^{N}\left(y_i - \hat{y}_i\right)^2
    \label{eq:loss}
\end{equation}

其中，$N$ 为样本数量，$y_i$ 为真实值，$\hat{y}_i$ 为预测值。

\section{图片示例}

图~\ref{fig:example} 展示了……

\begin{figure}[htbp]
    \centering
    % 替换为实际图片路径（相对于 figures/ 目录）
    % \includegraphics[width=0.6\textwidth]{example.pdf}
    \fbox{\parbox{0.5\textwidth}{\centering\vspace{2cm}图片占位符\vspace{2cm}}}
    \caption{示例图片标题}
    \label{fig:example}
\end{figure}

若需要中英文双语标题，可使用 \verb|\bilingualfigure| 命令：

\begin{figure}[htbp]
    % \includegraphics[width=0.5\textwidth]{example.pdf}
    \fbox{\parbox{0.5\textwidth}{\centering\vspace{1.5cm}双语图片占位符\vspace{1.5cm}}}
    \bilingualfigure{双语图标题（中文）}{Bilingual Figure Caption (English)}
    \label{fig:bilingual}
\end{figure}

\section{表格示例}

表~\ref{tab:example} 给出了……

\begin{table}[htbp]
    \centering
    \caption{示例三线表}
    \label{tab:example}
    \begin{tabular}{lccc}
        \toprule
        \textbf{方法} & \textbf{精度 (\%)} & \textbf{召回率 (\%)} & \textbf{F1 值} \\
        \midrule
        方法 A & 92.3 & 88.7 & 90.5 \\
        方法 B & 94.1 & 91.2 & 92.6 \\
        \textbf{本方法} & \textbf{95.8} & \textbf{93.4} & \textbf{94.6} \\
        \bottomrule
    \end{tabular}
\end{table}

\section{算法示例}

算法~\ref{alg:example} 描述了……

\begin{algorithm}[htbp]
    \caption{示例算法}
    \label{alg:example}
    \KwIn{输入数据 $\mathcal{D}$，迭代次数 $T$}
    \KwOut{模型参数 $\theta$}
    初始化参数 $\theta_0$\;
    \For{$t = 1$ \KwTo $T$}{
        计算梯度 $g_t \leftarrow \nabla_\theta \mathcal{L}(\theta_{t-1}; \mathcal{D})$\;
        更新参数 $\theta_t \leftarrow \theta_{t-1} - \eta \cdot g_t$\;
    }
    \Return $\theta_T$\;
\end{algorithm}

\section{代码示例}

\begin{lstlisting}[language=Python, caption={Python 代码示例}]
import numpy as np

def mean_squared_error(y_true, y_pred):
    """均方误差损失函数"""
    return np.mean((y_true - y_pred) ** 2)

# 示例调用
y_true = np.array([1.0, 2.0, 3.0])
y_pred = np.array([1.1, 1.9, 3.2])
loss = mean_squared_error(y_true, y_pred)
print(f"MSE Loss: {loss:.4f}")
\end{lstlisting}
